\documentclass[12pt]{article}
\usepackage[utf8]{inputenc}
\usepackage[T1]{fontenc}
\usepackage{hyperref}
\title{Udi Language Resources for K-12 Curriculum Development}
\author{Compiled by Claude Sonnet 5 (claude-sonnet-5) for Andr\'as Kornai}
\date{}
\begin{document}
\maketitle
\section{Introduction}
Udi (ISO 639-3 udi) is a Nakh-Daghestanian language. Its main community is the village of Nij in Qabala District, Azerbaijan. Sources give about 4,000 speakers in Nij and about 8,000 across several countries. This dossier rests on a single search round and is incomplete.
\section{Primary Educational Resources}
No specific Udi textbook or primer was found. Sources disagree on schooling. Omniglot says Udi has been taught in Nij primary schools since 1992. The Chaikhana article says children learn it through fourth grade and that some community members want it extended to eighth grade. An ECLING source says all formal teaching in Nij is in Azerbaijani, with no Udi classes. This conflict needs checking on the ground.
\section{Linguistic and Reference Materials}
\begin{itemize}
\item Wikipedia article on the Udi language: \url{https://en.wikipedia.org/wiki/Udi_language}
\item Omniglot page on Udi writing. It reports a shift from Cyrillic to Latin in the late 1990s: \url{https://www.omniglot.com/writing/udi.htm}
\item ECLING page: \url{https://titus.fkidg1.uni-frankfurt.de/ecling/ecling03.htm}
\end{itemize}
\section{Religious and Community Resources}
\begin{itemize}
\item Chaikhana article ``Udis: The Descendants of The Albanians'': \url{https://chaikhana.media/en/stories/9/udis-the-descendants-of-the-albanians}
\item Everyculture entry ``Udis - Orientation'': \url{https://www.everyculture.com/Russia-Eurasia-China/Udis-Orientation.html}
\end{itemize}
\section{Digital and Audio Resources}
No audio, video or corpus resources were found.
\section{Curriculum Development Context}
The Chaikhana article says the community has a 52-letter Latin-based alphabet. A curriculum would need to confirm the current status of Udi teaching in Nij, and to locate textbooks, possibly through Azerbaijani or Udi community channels. The search did not surface any.
\end{document}